% This file provides examples of some useful macros for typesetting
% dissertations.  None of the macros defined here are necessary beyond
% for the template documentation, so feel free to change, remove, and add
% your own definitions.
%
% We recommend that you define macros to separate the semantics
% of the things you write from how they are presented.  For example,
% you'll see definitions below for a macro \file{}: by using
% \file{} consistently in the text, we can change how filenames
% are typeset simply by changing the definition of \file{} in
% this file.
%
%% The following is a directive for TeXShop to indicate the main file
%%!TEX root = main.tex

\newcommand{\NA}{\textsc{n/a}}	% for "not applicable"
\newcommand{\eg}{e.g.,\ }	% proper form of examples (\eg a, b, c)
\newcommand{\ie}{i.e.,\ }	% proper form for that is (\ie a, b, c)
\newcommand{\etal}{\emph{et al}}

% Some useful macros for typesetting terms.
\newcommand{\file}[1]{\texttt{#1}}
\newcommand{\class}[1]{\texttt{#1}}
\newcommand{\latexpackage}[1]{\href{http://www.ctan.org/macros/latex/contrib/#1}{\texttt{#1}}}
\newcommand{\latexmiscpackage}[1]{\href{http://www.ctan.org/macros/latex/contrib/misc/#1.sty}{\texttt{#1}}}
\newcommand{\env}[1]{\texttt{#1}}
\newcommand{\BibTeX}{Bib\TeX}

% Define a command \doi{} to typeset a digital object identifier (DOI).
% Note: if the following definition raise an error, then you likely
% have an ancient version of url.sty.  Either find a more recent version
% (3.1 or later work fine) and simply copy it into this directory,  or
% comment out the following two lines and uncomment the third.
\DeclareUrlCommand\DOI{}
\newcommand{\doi}[1]{\href{http://dx.doi.org/#1}{\DOI{doi:#1}}}
%\newcommand{\doi}[1]{\href{http://dx.doi.org/#1}{doi:#1}}

% Useful macro to reference an online document with a hyperlink
% as well with the URL explicitly listed in a footnote
% #1: the URL
% #2: the anchoring text
\newcommand{\webref}[2]{\href{#1}{#2}\footnote{\url{#1}}}

% epigraph is a nice environment for typesetting quotations
\makeatletter
\newenvironment{epigraph}{%
	\begin{flushright}
	\begin{minipage}{\columnwidth-0.75in}
	\begin{flushright}
	\@ifundefined{singlespacing}{}{\singlespacing}%
    }{
	\end{flushright}
	\end{minipage}
	\end{flushright}}
\makeatother

% \FIXME{} is a useful macro for noting things needing to be changed.
% The following definition will also output a warning to the console
\newcommand{\FIXME}[1]{\typeout{**FIXME** #1}\textbf{[FIXME: #1]}}

% Author comment system: margin-style notes for draft feedback from
% you, your advisor, and collaborators.  Set showcomments to false
% before final submission to hide all notes at once.
\usepackage{ifthen}
\newboolean{showcomments}
\setboolean{showcomments}{true}
\ifthenelse{\boolean{showcomments}}
{
	\definecolor{myyellow}{RGB}{255, 228, 26}
	\definecolor{myblue}{RGB}{50, 50, 220}
	\newcommand{\nb}[2]{
		{\sf
			\fcolorbox{myyellow}{yellow}{\scriptsize\textbf{#1}}%
			$\blacktriangleright$%
			{\color{myblue}\fontsize{7pt}{8pt}\selectfont\textbf{#2}}%
		}%
	}
}
{
	\newcommand{\nb}[2]{}
}

% Define one shortcut per author/reviewer, e.g.:
%   \supervisor{Have you considered X?}  produces a highlighted note
% (note: \advisor is already taken by ubcdiss.cls)
\newcommand{\supervisor}[1]{\nb{Supervisor}{#1}}
\newcommand{\me}[1]{\nb{Me}{#1}}

% Tool/system names used across multiple chapters should be defined
% here (rather than in individual chapter files) to avoid "already
% defined" errors.  For example:
% \newcommand{\mytool}{\textsc{MyTool}\xspace}

% Research questions list environment (usable across multiple chapters)
\newlist{researchquestions}{enumerate}{1}
\setlist[researchquestions]{label*=\textbf{RQ\arabic*}}

% ===== CHAPTER-PREFIXED RESEARCH QUESTION NUMBERING SYSTEM =====
% RQ formatting that respects UBC thesis margins.
% Automatically uses the chapter number (RQ2.1 in Chapter 2, RQ3.1 in
% Chapter 3, etc.)
\newenvironment{rqlist}{%
    \begin{list}{}{%
    \setlength{\leftmargin}{0em}%
    \setlength{\labelwidth}{0em}%
    \setlength{\labelsep}{0em}%
    \setlength{\itemindent}{0em}%
    \setlength{\listparindent}{0em}%
    \setlength{\rightmargin}{0em}%
    \setlength{\itemsep}{0.75em}%
    \setlength{\parsep}{0em}%
    }%
  }{%
    \end{list}%
}

% Command for RQ items with automatic chapter numbering
% Usage: \rqitem{1}{Question text}
% In Chapter 2, produces: RQ2.1: Question text
\newcommand{\rqitem}[2]{%
    \item \textbf{RQ\thechapter.#1:} #2%
}
% If you prefer numbering offset by one (RQ1.x in Chapter 2, matching
% "research question 1 is answered in chapter 2"), use this instead:
% \newcommand{\rqitem}[2]{%
%     \item \textbf{RQ\number\numexpr\thechapter-1\relax.#1:} #2%
% }
% ===== END RESEARCH QUESTION NUMBERING SYSTEM =====


% Code formatting command (usable across multiple chapters)
\newcommand{\code}[1]{{\small\ttfamily\texttt{#1}}}

% END
